\documentclass[11pt]{article}
\usepackage[OT1]{fontenc}
\renewcommand{\rmdefault}{cmr}
\renewcommand{\sfdefault}{cmss}
\renewcommand{\ttdefault}{cmtt}
\usepackage{amsmath,amssymb,amsthm}
\usepackage[hidelinks]{hyperref}
\setlength{\textwidth}{6.1in}
\setlength{\oddsidemargin}{0.2in}
\setlength{\evensidemargin}{0.2in}
\setlength{\textheight}{9in}
\setlength{\topmargin}{-0.35in}
\setlength{\headheight}{13.6pt}
\setlength{\emergencystretch}{1.5em}
\newtheorem{theorem}{Theorem}[section]
\newtheorem{lemma}[theorem]{Lemma}
\newtheorem{proposition}[theorem]{Proposition}
\newtheorem{corollary}[theorem]{Corollary}
\theoremstyle{definition}
\newtheorem{definition}[theorem]{Definition}
\newtheorem{example}[theorem]{Example}
\newcommand{\Comp}{\operatorname{Comp}}
\newcommand{\Forb}{\operatorname{Forb}_{\rm ind}}
\newcommand{\E}{\mathsf E}
\newcommand{\C}{\mathcal C}
\newcommand{\B}{\mathcal B}
\newcommand{\OO}{\mathcal O}
\newcommand{\NN}{\operatorname{N}}
\newcommand{\dSigma}{d\text{-}\Sigma}

\title{Equality and universal comparison\\of positive exponential terms}
\author{}
\date{21 October 2026\\Working manuscript}
\begin{document}
\maketitle
\begin{abstract}
We study expressions generated from positive integer variables, the constant
$1$, addition, and $\E(f,g)=2^f g$. Universal strict and weak comparison
are $\Pi^0_1$-complete already at exponential depth two. In contrast,
extensional equality is decidable with any finite number of independent
variables, and universal comparison is decidable with one independent
variable at arbitrary finite depth. The proofs give an existential
definition of multiplication, recursive normal forms, and computable
eventual-sign bounds. We also transfer the obstruction to bounded
two-component vectors in Wilken--Weiermann's assignment syntax.
The results refute unrestricted effectiveness of the source's semantic
comparison, while leaving open effective restricted comparisons adequate
for actual derivation assignments.
\end{abstract}

\section{A comparison problem with two distinct boundaries}
An identity asserts that two expressions agree at every input. A
universal inequality asserts that their difference never has a
counterexample. For the expressions considered here these two tasks have
different computability behavior. The distinction persists although every
individual evaluation is a finite calculation on positive integers.

Let $\mathcal F_d$ be the finite terms generated by
\[
  1,\quad x_1,\ldots,x_d,\quad f+g,\quad
  \E(f,g)=2^f g .
\]
Variables range independently over $\mathbb N_{>0}$. A positive integer
constant is a sum of copies of $1$; zero is not a source term.
Exponential depth counts nesting of $\E$, with addition taking the
maximum of its children's depths. Write $f\prec g$ and $f\preceq g$ for
strict and weak inequalities holding at every positive input tuple.

For example, $\E(x+1,y)+\E(x,y)$ and three copies of $\E(x,y)$
have the same value $3y2^x$. Our normal form recognizes this identity
without testing inputs. On the other hand, a sum of squares of
differences can vanish exactly at the solutions of an arbitrary
Diophantine equation. We encode that test using the allowed syntax,
with no subtraction, zero, or ordinary variable multiplication in the
constructed terms.

There are three main conclusions:
\begin{enumerate}
\item universal comparison with an unrestricted finite collection of
independent inputs is undecidable even at depth two;
\item equality is decidable uniformly in that finite collection and
at any finite depth;
\item with one independent input, strict and weak universal comparison
are decidable at any finite depth, and a false comparison has a
computable least counterexample.
\end{enumerate}
The number of variables in the undecidability reduction grows with the
input. We do not establish a smallest fixed number of variables at which
comparison becomes undecidable.

\subsection{The source question}
Wilken and Weiermann define semantic comparison by quantifying over
admissible substitutions and conjecture its effectiveness immediately
after their Definition~2.8 \cite{WW}. Their definitions allow the fragment
above: ground-component substitutions take positive finite values, and
every positive integer has a closed representation as a sum of $1$'s.
Natural sum and natural product are ordinary addition and multiplication
on these values. Their ordinal norm $\NN(n)$ equals $n$ for finite $n$.
Thus their comparison on these terms is exactly comparison over
$\mathbb N_{>0}^d$, not an approximation to it.

The variable multiplier in $\E(f,g)$ matters. We are not studying only
addition and the unary operation $f\mapsto2^f$. The use of \emph{unary}
here counts independent numerical inputs.

The undecidability proof uses MRDP as its published input. The identity
and one-variable procedures use the elementary separation argument below,
rather than a transcendence assumption. The final vector result uses
the source's own norm and collapsing-function estimates.

\section{Universal comparison at depth two}
\begin{theorem}[Finite depth-two obstruction]
\label{thm:finite-bare-hardness}
For either $\triangleleft\in\{<,\le\}$, consider pairs $(A,B)$ in the
finite fragment $\mathcal F_d$, with $d$ allowed to vary, satisfying
\[
 \forall\chi\text{ admissible in Definition~2.7 of \cite{WW}},\qquad
 \operatorname{val}(A\chi)\mathrel{\triangleleft}
 \operatorname{val}(B\chi)
\]
is $\Pi^0_1$-complete already when the terms use only $1$, independently
named ground-type variables, $+$ and $\E$, and have $\E$-nesting depth
at most two. Every value is a positive integer, and each output term
belongs to every $\C_i$, $i\ge1$. The same restricted completeness
holds for comparison of their norms.
\end{theorem}

For these variables, the original admissible domain is exactly the set
of positive integer tuples. Indeed closed $\B_0$-terms are positive
finite, every positive integer is represented by a sum of $1$'s, and
the bounded-norm condition on a one-component substitution is automatic
\cite[Definitions~2.5--2.7]{WW}. All variables introduced below are
fresh ground-type variables of this kind.

\subsection{A bounded multiplication identity}

\begin{lemma}\label{lem:restricted-product}
Let $a,y,s$ be positive integers and suppose
$y<2^a-1$ and $s<a(2^a-1)$. Then $s=ay$ if and only if there exist
positive integers $u,q$ such that
\begin{align}
 2^a u+1&=u+2^s,\label{eq:mersenne-first}\\
 2^a q+y+1&=u+q+2^a.\label{eq:mersenne-second}
\end{align}
\end{lemma}

\begin{proof}
Write $B=2^a$ and $D=B-1$. Equation~\eqref{eq:mersenne-first} says
$Du=2^s-1$. Division of $s$ by $a$, say $s=ka+r$ with $0\le r<a$,
gives $2^s-1\equiv2^r-1\pmod D$. Since
$0\le2^r-1<D$, divisibility forces $r=0$. Thus $k=s/a\ge1$ and
\[
 u=1+B+\cdots+B^{k-1}\equiv k\pmod D.
\]
Equation~\eqref{eq:mersenne-second} is $D(q-1)=u-y$, so
$k\equiv y\pmod D$. The two bounds give $1\le k,y<D$, whence
$k=y$ and $s=ay$.

Conversely, if $s=ay$, take
\[
 u=\sum_{j=0}^{y-1}B^j,
 \qquad q=1+\frac{u-y}{D}.
\]
The second expression is an integer because $u\equiv y\pmod D$,
and it is positive because $u\ge y$. Both equations follow.
This includes $y=1$, when $u=q=1$.
\end{proof}

The lemma's bounds and equations are expressible using $1,+,\E$:
they contain no product of two arbitrary variables. The apparent
products $2^a u$, $2^a q$ and $a2^a$ are exactly permitted
exponential-constructor terms.

\begin{lemma}[Seven-equation definition of multiplication]
\label{lem:seven-equations}
For positive $x,y,w$, put $a=x+y+1$ and $b=a+x$, treating these as
affine term abbreviations. Then $w=xy$ if and only if the following
seven source equations have a solution in positive integers
$s,t,u,v,q,r,h,k$:
\begin{align*}
 \E(a,u)+1&=u+\E(s,1),&
 \E(b,v)+1&=v+\E(t,1),\\
 \E(a,q)+y+1&=u+q+\E(a,1),&
 \E(b,r)+y+1&=v+r+\E(b,1),\\
 s+a+h&=\E(a,a),&
 t+b+k&=\E(b,b),\\
 t&=s+w.&&
\end{align*}
Every side has $\E$-nesting depth at most one.
\end{lemma}

\begin{proof}
Since $a\ge y+2$ and $b>a$, the bounds
$y<2^a-1$ and $y<2^b-1$ hold identically. The positive slacks $h,k$
give $s<a(2^a-1)$ and $t<b(2^b-1)$. Lemma~\ref{lem:restricted-product}
therefore forces $s=ay$ and $t=by$. The last equation now gives
$w=(b-a)y=xy$.

Conversely, if $w=xy$, choose $s=ay$ and $t=by$; use the witnesses
from the preceding lemma for the two pairs of equations, and set
\[
 h=a(2^a-1-y)>0,\qquad k=b(2^b-1-y)>0.
\]
All seven equations hold. Products in these witness formulas are
ordinary mathematical calculations; the displayed source equations
contain only $1,+,\E$. No disjunction, zero-valued variable or
extra quantifier alternation is needed.
\end{proof}

\subsection{From Diophantine equations to one universal comparison}

\begin{proof}[Proof of Theorem~\ref{thm:finite-bare-hardness}]
Solvability of integer-coefficient polynomial equations over positive
integers is $\Sigma^0_1$-complete by the MRDP theorem
\cite{Matiyasevich}. The positive convention follows from the usual
nonnegative convention by replacing each nonnegative variable with
$x-1$ and expanding the polynomial.

Given such a polynomial $P$, separate its positive and negative
coefficient parts and express
\[
 P_+(\mathbf x)+1=P_-(\mathbf x)+1.
\]
Each side has a finite arithmetic circuit using positive constants,
addition and multiplication. Delete zero monomials; the added $1$
ensures that both sides are positive, even if one part was empty.
Introduce a fresh positive variable for every circuit wire. Addition
and constant gates give affine equations. Replace each multiplication
gate by Lemma~\ref{lem:seven-equations}, using fresh auxiliaries.
Retain the final output equality. We obtain a finite system
\[
 f_i(\mathbf z)=g_i(\mathbf z),\qquad 1\le i\le m,
\]
over positive variables which has a solution exactly when $P$ does.
All its terms use only $1,+,\E$, and every $f_i,g_i$ has
$\E$-depth at most one. This is an effective syntactic construction;
the new variables, like the original ones, become independently
substitutable ground-type variables.

Form the source terms
\begin{align*}
 U&=1+\sum_{i=1}^{m}
       \bigl(\E(f_i+f_i,1)+\E(g_i+g_i,1)\bigr),\\
 V&=1+\sum_{i=1}^{m}\E(f_i+g_i+1,1).
\end{align*}
At every positive tuple, ordinary integer arithmetic gives
\begin{equation}\label{eq:finite-gap}
 U-V=\sum_{i=1}^{m}\left(2^{f_i}-2^{g_i}\right)^2\ge0.
\end{equation}
The right side is a proof identity, not source syntax. It vanishes
exactly when all $f_i=g_i$. Consequently
\[
 P\text{ has no positive zero}
 \quad\Longleftrightarrow\quad V\prec U
 \quad\Longleftrightarrow\quad V+1\preceq U.
\]
Here the final equivalence uses that the nonnegative difference is an
integer. These are computable many-one reductions from a
$\Pi^0_1$-complete problem, with output depth at most two.

For the upper bound on this restricted family, enumerate all positive
integer substitutions and evaluate the finite terms exactly. A failed
comparison is witnessed by one enumerated tuple, so the relation is
co-c.e. This establishes the claimed upper bound on the finite fragment.

For every $i\ge1$, simultaneous induction on these finite source terms
gives membership in $\C_i$: ground variables lift from $\C_0$, sums
preserve membership, and $\E(f,g)$ uses $f\in\C_{i+1}$ and
$g\in\C_i$ \cite[Definition~2.9]{WW}. Finally all values are finite,
so $\NN(U)=U$ and $\NN(V)=V$. The identical reduction therefore
proves the asserted norm-comparison completeness.
\end{proof}

\subsection{\texorpdfstring{The obstruction persists on bounded $\C$-vectors}{The obstruction persists on bounded C-vectors}}

\begin{theorem}[Two-component bounded-vector obstruction]
\label{thm:bounded-bare-vectors}
Both strict and weak comparison in Definition~2.8 are undecidable even
on a computably constructed family of two-component $\C$-vectors whose
bounded norm is proved uniformly. The family gives $\Pi^0_1$-hardness
while retaining the finite depth-two terms as higher components.
\end{theorem}

\begin{proof}
Use the $U,V$ just constructed and put $R=U+V+1$. Let $L\ge1$ be
the number of nodes in the source term tree of $R$, and let
$z_1,\ldots,z_n$ list all its variables. Set
\[
 s_0=1+\sum_{j=1}^{n}z_j,
 \qquad b=(L+1)s_0,
 \qquad g=5b.
\]
These are positive affine source terms; the integer coefficients
abbreviate finite sums. We claim $R\preceq F_2(g)$.

Fix a substitution. Every leaf value is at most $s_0$ and every
depth-zero subtree has value at most $b$. Flatten sums at each depth.
There are at most $L$ summands, so every depth-one subtree is bounded by
\[
 Lb2^b\le2^{3b}.
\]
At depth two this gives
\[
 R\le L\,2^{\,2^{3b}+3b}
   \le2^{\,2^{3b}+4b}
   \le2^{\,2^{5b}}.
\]
We used $L,b\le2^b$ and $2^{3b}+4b\le2^{5b}$ for $b\ge1$.
The source defines $F_0(n)=2^n$ and
$F_{j+1}(n)=F_j^{\,n+1}(n)$, so
$F_2(g)\ge F_1(g)\ge2^{2^g}$ \cite[p.~15]{WW}.
This proves the claim over exactly the admissible substitution domain.

Now define the source term
\[
 H=\psi(\omega\mathbin{\cdot}R+g).
\]
We have $R\in\C_1$, $g\in\C_0$ and
$\NN(R)=R\preceq F_2(g)$, which is precisely the side condition
for $H\in\C_0$. Under each substitution, natural product gives
$\omega\otimes R=\omega\cdot R$, and hence
\[
 H\ge\NN(\omega\otimes R\oplus g)=2R+g\ge R
\]
by \cite[Proposition~2.3]{WW}. Therefore
$(H,V+1)$, $(H,U)$ and $(H+1,U)$ are bounded $\C$-vectors.
Definition~2.8 gives exactly
\[
 (H,V+1)\preceq(H,U)
 \quad\Longleftrightarrow\quad V+1\preceq U
 \quad\Longleftrightarrow\quad
 (H,V+1)\prec(H+1,U),
\]
because $H\prec H+1$ holds under every substitution. Applying
\eqref{eq:finite-gap} proves the reduction.
\end{proof}

This result directly concerns the original scalar and vector comparison
relations, including bounded $\C$ inputs. It does not establish that the
hard vectors occur as assignments of actual $T$-terms, and it does not
rule out effective sufficient comparisons or a differently designed
termination ordering. It rules out an algorithm deciding all the stated
semantic comparisons exactly.

\section{Normal forms and extensional equality}
The auxiliary normal-form syntax permits zero.  Define, by finite
well-founded trees, expressions
\[
             F=\sum_{e\in\operatorname{supp}(F)}P_e(\mathbf x)2^e,
             \qquad P_e=c_e+\sum_{i=1}^d a_{e,i}x_i,
             \tag{N1}\label{eq:normal-form}
\]
where every $P_e$ is a nonzero affine polynomial with nonnegative integer
coefficients, the support is finite, and each key $e$ is itself a normal
form of smaller height.  Every key is \emph{constant-free}: its
coefficient of $1$ in the zero-exponent term is zero.  Zero is represented
by the empty sum.  Equal keys are merged, and a zero coefficient is
deleted.  A normal form is a finite map, so the display order of its keys
has no mathematical significance.

For any normal form $F$, let $c(F)$ denote the constant coefficient of
its zero-exponent term, with value zero when that term is absent.  Let
$F^{\circ}$ be the form obtained by removing just that coefficient.
Thus $F=c(F)+F^{\circ}$ and $F^{\circ}$ is constant-free.  The name
``constant-free'' is syntactic: for example, $2^{x_1}$ has $c(F)=0$
although its value at $x_1=0$ would be $1$.

\begin{lemma}[Effective normalization]\label{lem:normalization}
Every term in $\mathcal F_d$ has an effectively computable normal form
$\operatorname{nf}(f)$ of the shape \eqref{eq:normal-form}, with the same value
at every positive integer tuple.  All operations terminate by finite
syntax recursion.
\end{lemma}
\begin{proof}
The normal forms of $1$ and $x_i$ have only the zero key.  Addition merges
the two finite maps and adds their affine coefficients.  Suppose $F,G$
are the normal forms of $f,g$, and write $F=c+R$, with $R$ constant-free.
Then use
\[
     2^F G
       =\sum_{e\in\operatorname{supp}(G)}
             (2^cP_e)\,2^{R+e}.                         \tag{N2}
\]
Addition of the two already normalized constant-free forms $R,e$ gives
a constant-free key.  Merge equal resulting keys.  The coefficients
remain affine because only addition and multiplication by a numerical
constant occurred.  The identities used are ordinary integer
distributivity and $2^{u+v}=2^u2^v$.  Keys remain of smaller finite height
than the resulting outer sum.  Computing $2^c$ may be expensive, but it
is a finite computation; no complexity bound is being asserted.
\end{proof}

\subsection{A separation lemma without transcendence assumptions}

\begin{lemma}[Finite-family ray separation]\label{lem:ray}
Given a finite collection $\mathcal S$ of normal forms, close it under
taking keys and add zero.  One can effectively choose positive integers
$a_1,\ldots,a_d$ such that, on the ray $x_i=a_it$ with integer $t>0$,
the following statements hold for all $F,G\in\mathcal S$.
\begin{enumerate}
\item If $F\ne G$, their difference is eventually nonzero with a fixed
      sign.
\item If $F,G$ are constant-free and distinct, then
      $|F(\mathbf a t)-G(\mathbf a t)|\ge t$ for all sufficiently large
      $t$.
\end{enumerate}
The common thresholds and the eventual signs are computable.
\end{lemma}

\begin{proof}
Include the zero affine coefficient when collecting the affine
coefficient vectors of all forms in $\mathcal S$.  Choose a positive
integer $B$ so large that, for every difference $v$ of two collected
vectors of variable coefficients,
\[
             B>1+\sum_i |v_i|.
\]
Set $a_i=B^{i-1}$.  If $v\ne0$ and $j$ is its highest nonzero index,
the absolute value of $v_ja_j$ exceeds the sum of the absolute values of
the earlier contributions.  Thus $v\cdot\mathbf a\ne0$.
Consequently every nonzero difference of two collected affine
coefficients becomes a nonzero affine polynomial in $t$.  A difference
with nonzero variable part acquires a nonzero integer slope.

Induct on normal-form height, simultaneously for pairs in the finite
collection.  When only the zero key occurs, there is an affine polynomial
$At+C$.  If $A\ne0$, its eventual sign is the sign of $A$, and
$t\ge |C|+1$ suffices to make it nonzero.  If $A=0$, its nonzero constant
sign is immediate.  When both forms are constant-free, $C=0$, so a
nonzero difference has absolute value at least $t$.

For the induction step, subtract the two maps and discard zero
coefficients.  Distinct remaining keys are constant-free forms of
smaller height.  By induction their pairwise eventual order is strict,
and any ordered pair has a gap at least $t$ after a computable threshold.
Let $e_*$ be the greatest remaining key in that eventual order, with
coefficient $P_*(t)=A_*t+C_*$.  For all sufficiently large $t$,
$P_*(t)$ has a fixed sign and integer absolute value at least one.
Put
\[
       C=\sum_{e\ne e_*} (|A_e|+|C_e|).
\]
Once every key gap $e_*(t)-e(t)$ is at least $t$, the sum of the absolute
values of the lower contributions is at most
\[
        C t\,2^{e_*(t)-t}.                              \tag{N3}
\]
For $t\ge\max(4,2C)$, the elementary inequality $2^t\ge t^2$
implies $Ct2^{-t}\le1/2$.  Thus the leading summand determines the
sign of the difference.

If $e_*=0$, there were no other surviving keys, since every nonzero
constant-free key is eventually positive.  This is the already treated
affine case.  If $e_*>0$, induction applied to
$e_*$ and zero gives $e_*(t)\ge t$ after another computable threshold.
The difference then has absolute value at least
$2^{e_*(t)-1}\ge t$.  This proves both assertions.  Taking maxima of
the finitely many recursive thresholds makes every step effective.
No algebraic-independence or transcendence theorem is used.
\end{proof}

\begin{theorem}[Decidable extensional equality]\label{thm:identity}
For any finite number of independent positive ground variables and any
finite exponential nesting depths,
\[
  (\forall\mathbf n>0)\ f(\mathbf n)=g(\mathbf n)
       \quad\Longleftrightarrow\quad
  \operatorname{nf}(f)=\operatorname{nf}(g).                 \tag{N4}
\]
Consequently extensional equality in $\mathcal F_d$ is decidable,
uniformly in $d$ and the two input terms.  If the normal forms differ,
a separating positive integer tuple can be computed.
\end{theorem}
\begin{proof}
Equal normal forms imply equality by normalization soundness.  For
unequal forms, apply Lemma~\ref{lem:ray} to the two forms and their
keys.  Their values differ on every sufficiently large point of the
computed positive ray.  The constructive proof supplies such a point.
Structural equality of the finite normal forms is decidable.
\end{proof}

\subsection{A complete universal comparison procedure with one variable}

\begin{theorem}[One-variable universal comparison]\label{thm:unary}
Let $f,g\in\mathcal F_1$.  Both
\[
                 (\forall n>0)\ f(n)<g(n),\qquad
                 (\forall n>0)\ f(n)\le g(n)                \tag{N5}
\]
are decidable, uniformly at arbitrary finite exponential depth.
A false instance has an effectively computable least positive
counterexample.
\end{theorem}

\begin{proof}
Normalize $f,g$.  Equality is handled immediately.  With one variable
the ray $x=t$ already separates every nonzero affine coefficient
difference, so the proof of Lemma~\ref{lem:ray} applies with $a_1=1$.
It computes a sign $s\in\{-1,1\}$ and integer $N\ge1$ such that
$\operatorname{sign}(f(n)-g(n))=s$ for all $n\ge N$.

For clarity, one permissible threshold recurrence is explicit.  For a
nonzero affine difference $An+C$, take $N=1$ if $A=0$, and
$N=|C|+1$ otherwise.  For a difference whose greatest surviving key
$e_*$ is nonzero, take the maximum of:
\begin{enumerate}
\item $4$ and $2\sum_{e\ne e_*}(|A_e|+|C_e|)$;
\item $|C_*|+1$ if $A_*\ne0$, and $1$ otherwise;
\item the recursively obtained thresholds at which
      $e_*(n)-e(n)\ge n$ for every other surviving key $e$, and
      $e_*(n)\ge n$.
\end{enumerate}
Keys are constant-free, so their distinct-pair recursive differences
always have the gap property required in the third item.  The estimates in the ray-separation proof verify this recurrence.  It is a recursion
on strictly smaller key heights, not a search for an unspecified
asymptotic bound.

Evaluate the two terms exactly at $1,\ldots,N-1$.  If one of these is
a counterexample, return the first.  If none is, the eventual sign
either proves the comparison for all remaining integers or makes $N$
its least counterexample.  Every integer evaluation terminates.
This proves a total computability theorem, without claiming a practical
bound on the magnitudes of the evaluated integers.
\end{proof}

\begin{corollary}[Further exact consumers]\label{cor:consumers}
The same procedures decide equality of finite-valued ordinal vectors
over any finite set of independent ground variables, and strict/weak
vector comparison and bounded norm for finite-valued $\mathcal O$-vectors
with one independent ground variable.  They also decide comparison
after any fixed affine one-parameter substitution with nonnegative
integer coefficients and positive coordinate values for finitely many
ground variables: explicitly,
$x_i=a_i n+b_i$ with $a_i,b_i\in\mathbb N$ and $a_i+b_i>0$.
\end{corollary}
\begin{proof}
Vector equality and comparisons are the componentwise combinations
specified by Wilken--Weiermann Definition~2.8.  On finite values the
norm equals the value.  Such an affine substitution is a term built
from the remaining variable, positive integer constants and addition;
substitution therefore reduces its comparison to
Theorem~\ref{thm:unary}.  A zero constant summand is omitted.  A
coordinate that is a positive constant is allowed as well.
\end{proof}


\section{What the algorithms do and do not decide}
The normal form is an identity test, not a universal multivariable
ordering: a ray witnessing a nonidentity does not show that an inequality
holds away from that ray. This is why the identity algorithm is
consistent with the depth-two lower bound. With one variable there
are only a finite initial segment and the eventual tail, so the same
asymptotic argument becomes a complete comparison procedure.

A fixed substitution $x_i=a_i n+b_i$ with nonnegative integer
coefficients and $a_i+b_i>0$ also reduces comparison to the unary case.
Zero coefficients are handled by omitting summands; a positive constant
coordinate is permitted. This is a restriction to a prescribed path
through the input domain, not a decision for independent variables.

The explicit threshold calculations can produce enormous integers.
They prove total computability, not a useful running-time estimate.
The bounded-vector reduction likewise proves hardness for the specified
source syntax, without asserting that the hard vectors are realized by
assignments of actual terms of G\"odel's $T$. An algorithm for a justified
assignment image, or a sound incomplete comparison adequate for
termination proofs, is a distinct target. None of the results constructs
an exact transitive termination ordering with a matching ordinal height.

\paragraph{Preparation and status.}
This manuscript was prepared with AI assistance from a research programme
directed by the repository owner. The displayed unbounded arguments,
rather than numerical trials, support the theorems. MRDP and the
specified Wilken--Weiermann statements are cited inputs. No independent
human review, proof-assistant verification, or historical priority claim
is asserted. The canonical manuscript and build files are included;
no private proof notes or project conventions are required.

\begin{thebibliography}{9}
\bibitem{WW} G.~Wilken and A.~Weiermann,
\emph{Derivation Lengths Classification of G\"odel's $T$ Extending
Howard's Assignment}, Logical Methods in Computer Science 8(1:19)
(2012), 1--44. Source version arXiv:1203.0115v2,
Definitions~2.4--2.9 and Proposition~2.3.
\url{https://doi.org/10.2168/LMCS-8(1:19)2012}.
\bibitem{Matiyasevich} Y.~V.~Matiyasevich,
\emph{Hilbert's Tenth Problem}, MIT Press, 1993.
The MRDP theorem is used for Diophantine solvability; the reduction
from that problem to exponential comparison is proved here.
\end{thebibliography}
\end{document}
